\documentclass[11pt]{scrartcl}
\usepackage[sexy]{evan} %BLBLBLBLBLBLBLBLBLBLBLB EVANCHEN
\usepackage[T1]{fontenc}
\usepackage[utf8]{inputenc}
\usepackage{graphicx}
\usepackage{amsmath}
\usepackage{tikz}
\usepackage{tikz-cd}
\usepackage{asymptote}
\usepackage{amsthm}
\usepackage{amssymb}
\usepackage{gensymb}
\usepackage{amsfonts}
\usepackage[english,italian]{babel}
\usepackage{quoting}
\quotingsetup{font=big}
\usepackage{booktabs}
\usepackage{caption}
\usepackage{lmodern}
\usepackage{faktor}
\usepackage{bbm} %oer fare gli insiemi indicatori
\usepackage{leftindex} %per gli ideali e gli indici a sinistra
\usepackage{hyperref}
\usepackage{epigraph} 
\usepackage{pgfplots}
\usepackage[export]{adjustbox}

%Pacchetti per i codici
\usepackage{listings}
\usepackage{xcolor}
\usepackage{algpseudocode}

%Analisi e probabilità
%----------------------------------
\newcommand{\dif}{\text{d}}
%%% the pure symbol, see https://tex.stackexchange.com/a/718194/4427
\newcommand{\cind}{\perp\mspace{-9mu}\perp}
%%% define \indep
\NewDocumentCommand{\indep}{e{_}}{%
	\mathrel{%
		\IfNoValueTF{#1}{% no subscript
			\cind
		}{% subscript
			\mathchoice{\mathop{\cind}\limits_{#1}}{\cind_{#1}}{\cind_{#1}}{\cind_{#1}}%
		}%
	}%
}
%-----------------------------------

%Algebra lineare e geometria
%----------------------------------
\newcommand{\tr}[1]{\text{tr}\!\left(#1\right)}
\renewcommand{\rank}[1]{\text{rank}\!\left(#1\right)}
\newcommand{\Immagine}[1]{\text{Im}\left(#1\right)}
\renewcommand{\ker}[1]{\text{Ker}\left(#1\right)}
\renewcommand{\AA}{\mathbb{A}}
\newcommand{\LL}{\mathbb{L}}
\newcommand{\SO}{\text{SO}}
%----------------------------------

%Fisica
%----------------------------------
\newcommand{\UDiMisura}[1]{\,\text{#1}}
\newcommand{\ihat}{\hat{\imath}}
\newcommand{\jhat}{\hat{\jmath}}
\newcommand{\khat}{\hat{k}}
\newcommand{\posit}{\vec{r}}
\newcommand{\veloc}{\dot{\vec{r}}}
\newcommand{\accel}{\ddot{\vec{r}}}
%----------------------------------

%Informatica
%----------------------------------
\definecolor{codegreen}{rgb}{0,0.6,0}
\definecolor{codegray}{rgb}{0.5,0.5,0.5}
\definecolor{codepurple}{rgb}{0.58,0,0.82}
\definecolor{backcolour}{rgb}{0.95,0.95,0.92}

\lstdefinestyle{overleafwiki}{
	backgroundcolor=\color{backcolour},   
	commentstyle=\color{codegreen},
	keywordstyle=\color{magenta},
	numberstyle=\tiny\color{codegray},
	stringstyle=\color{codepurple},
	basicstyle=\ttfamily\footnotesize,
	breakatwhitespace=false,         
	breaklines=true,                 
	captionpos=b,                    
	keepspaces=true,                 
	numbers=left,                    
	numbersep=5pt,                  
	showspaces=false,                
	showstringspaces=false,
	showtabs=false,                  
	tabsize=2
}

\lstset{style=overleafwiki}
\newcommand{\bigo}[1]{\mathcal{O}\!\left(#1\right)}
\newcommand{\bigomega}[1]{\Omega\!\left(#1\right)}
\newcommand{\bigtheta}[1]{\Theta\!\left(#1\right)}
%----------------------------------

%Insiemistica e Algebra
%----------------------------------
\newcommand{\eset}{\varnothing}
\newcommand{\parti}[1]{\mathcal{P}\!\left(#1\right)}
\newcommand{\inj}{\hookrightarrow}
\newcommand{\surj}{\twoheadrightarrow}
\newcommand{\bij}{\hookrightarrow\mathrel{\mspace{-15mu}}\rightarrow}
\newcommand{\ordine}[2]{\text{ord}_{#1}\left(#2\right)}
\newcommand{\normal}{\mathrel{\unlhd}}
\newcommand{\MCD}{\text{MCD}}
\renewcommand{\mcm}{\text{mcm}}
%----------------------------------


\begin{document}
	\title{Controesempio per un Greedy}
	\subtitle{Matrix Chain Multiplication}
	\date{13 maggio 2026}
	
	\maketitle
	
	Definiamo due approcci Greedy per il problema della Matrix Chain Multiplication:
	\begin{description}
		\item[Top-Down:] Costruiamo l'albero binario completo associato alla parentesizzazione della catena di matrici in modo greedy dall'alto: per ogni istanza scegliamo le due matrici adiacenti con dimensione comune minima 
		$$A_{i}\in \CC_{a\times m}$$
		$$A_{i+1}\in \CC_{m\times b}$$ 
		dove $m$ è il minimo tra tutte le dimensioni comuni, e ricorriamo sui sottoproblemi $\langle 1,2,\dots,i\rangle$ e $\langle i+1,i+2,\dots,n\rangle$ (dove $n$ è il numero totale di matrici).
		\item[Bottom-Up:] Costruiamo l'albero binario completo associato alla parentesizzazione della catena di matrici in modo greedy dal basso: in modo analogo a come fatto a lezione per i codici di Huffmann, consideriamo inizialmente tutti i nodi scollegati e colleghiamo la coppia di nodi con dimensione comune maggiore, ovvero
		$$A\in \CC_{a\times M}$$
		$$B\in\CC_{M\times b}$$
		(dove $M$ è il massimo tra tutte le dimensioni comuni) e al loro posto sostituiamo la matrice $C\in \CC_{a\times b}$, ricorrendo in coda.
	\end{description}
	Per entrambi questi approcci un controesempio che ne dimostra la scorrettezza è il seguente: 
	\begin{example*}
		Vogliamo trovare la parentesizzazione ottimale per 
		\begin{gather*}
			A_1\in \CC_{1\times (M-1)} \\
			A_2\in \CC_{(M-1)\times M}\\
			A_3\in \CC_{M\times (M-1)}\\
			A_4\in \CC_{(M-1)\times 1}
		\end{gather*}
		entrambi gli algoritmi producono una delle due parentesizzazioni
		$$(A_1(A_2A_3))A_4$$
		$$A_1((A_2A_3)A_4)$$
		che hanno un costo computazionale uguale e pari a $M^3-M^2$. Tuttavia la parentesizzazione $(A_1A_2)(A_3A_4)$ ha un costo pari a $M^2-M$, e per $M$ sufficientemente grande si vede facilmente che
		$$M^3-M^2>M^2-M$$
		il che invalida gli algoritmi proposti sopra.
	\end{example*}
	
	
\end{document}
